\section{Consequences and boundaries}
\begin{corollary}\label{cor:basis13}
The set \(C_{12}\cup\{0\}\) is an asymptotic basis of exact order
thirteen: every sufficiently large integer is a sum of at most
thirteen elements of \(C_{12}\), and twelve do not suffice.
\end{corollary}
\begin{proof}
For a large target choose \(j\in\{0,\ldots,11\}\) with
\(N-j\equiv2\pmod{12}\). Theorem~\ref{thm:c12} writes \(N-j\)
as two core values; add \(j\) copies of one. This uses at most thirteen
positive summands. Conversely, \(25+60r\), \(r\ge0\), is
\(1\pmod{12}\) and divisible by the forbidden prime five.
At most twelve positive core summands would force exactly one
summand on that residue class, which is impossible. Padding with
zeros gives the equivalent fixed-order assertion.
\end{proof}

The field formulation follows from classical abelian class field
theory. A subgroup condition on Frobenius is complete splitting in
its fixed field. Finite conjunctions become intersections of inverse
image subgroups in the compositum. Kronecker--Weber places this
compositum in a cyclotomic field
\cite[Chapter I, Theorem 4.16]{milnecft}, giving precisely the
fixed subgroups in Theorem~\ref{thm:general}, after discarding the
fixed bad-prime set. The reduction is classical; the pointwise
indicator correlation is the analytic assertion used here.

These conclusions have several precise boundaries. The semigroup
theorems treat subgroup pairs and unfiltered self-sums for a prime
coset. Other unions of residue classes require their own model and
sieve checks. The modulus and coefficients stay fixed. The interval discrepancy is absolute,
and gives no relative asymptotic in intervals of power-short length
\(N^\theta\), \(\theta<1\). The threshold is not numerical, so it
does not classify small exceptions. The exact sumsets of primitive
quadratic norms and their finite certificates form a separate
arithmetic problem, treated in the companion article
\emph{Exact Sumsets of Primitive Quadratic Norms}.
The separate article
\href{https://hexagonmath.org/2611.00593v1}{\emph{Pointwise representation
counts for Eisenstein restricted sums}} addresses where summands may
be placed, with a different proof and parameter range.

The shared comparison requires a common invariant prime class.
Two arbitrary prime prescriptions can be disjoint, so that no such
class is available. Their binary correlation is not a consequence of
Theorem~\ref{thm:model-comparison}. Nor does its fixed-shift clause
assert uniformity for growing sublinear shifts between its two
proved frequency regimes.

Higher-fold sums and independent coupled systems can use classical
linear-form correlation theorems, but their proof structure differs
from this binary reflected line. The ordinary quadratic-form
indicator correlations in Section~\ref{sec:quadratic-pairs} use
the fixed-shift clause after separately verifying its arithmetic
hypotheses. They include the Connors--Keating case and Smilansky's
five-form family. Growing-shift uniformity and higher indicator
correlations remain outside those statements.
